\section{Conclusion}

This work presented a reproducible benchmarking methodology to evaluate reconstruction fidelity, adaptation stability, and computational efficiency across generalizable and per-scene neural rendering architectures under sparse-view conditions. In the experimental evaluation, per-scene models achieved the best visual quality and rendering speed under $N=10$ views: Nerfacto reached $27.11$ dB PSNR and $0.171$ LPIPS, while Splatfacto (3DGS) rendered novel views at $200.7$ FPS using $0.10$ GB of dynamic VRAM. Under three input views ($N=3$), modular fine-tuning of pre-trained priors prevented representational collapse: adapting GNT via \textit{PFT-Decoder} ($15.07$ dB) or \textit{PEFT-LoRA} ($14.57$ dB) reduced adaptation time by $16.3\%$ compared to full fine-tuning and outperformed Instant-NGP ($14.06$ dB), although limited by volumetric ray-marching throughput ($\le 0.015$ FPS).

Future work will extend this evaluation protocol to unbounded $360^\circ$ scenes and investigate depth-regularized sparse-view 3D Gaussian Splatting methods \cite{seo2026improvingsparseview3dgsgeneralization, song2025d2gsdepthanddensityguidedgaussian} alongside feed-forward Gaussian architectures.
